%% kapa: exam=1 round=28 year=2017 period=1 form=A subject=real_estate subject_ko=부동산학원론 seq=10 number=90
%% source: Q-Net 감정평가사 기출문제 (cst00309) artlSeq=5199603 file=extracted/1008/5199603/2017년 제28회 감정평가사 1차시험 문제지/1교시 A형.pdf
%% conversion: kordoc 4.24.51 -> tools/gen1.py (auto); review=proofread-sonnet; flags=none
\begin{question}{90}
부동산투자의 수익률에 관한 설명으로 옳지 않은 것은? (단, 주어진 조건에 한함)
\begin{choices}
  \item 기대수익률은 투자한 부동산의 예상수입과 예상지출로 계산되는 수익률이다.
  \item 실현수익률이란 투자가 이루어지고 난 후에 실제로 달성된 수익률이다.
  \item 요구수익률은 투자자에게 충족되어야 할 최소한의 수익률이다.
  \item 장래 기대되는 수익의 흐름이 주어졌을 때, 요구수익률이 클수록 부동산의 가치는 증가한다.
  \item 투자자의 요구수익률은 체계적위험이 증대됨에 따라 상승한다.
\end{choices}
\end{question}
\answer{4}
